A equação (2.12) do livro-texto foi derivada usando o Método Básico da Análise Dimensional para o problema da queda livre no vácuo, resultando na seguinte forma funcional [1]:
$$ V = \text{constante} \times \sqrt{gh} $$

\begin{itemize}
    \item \textbf{Qual é a constante?} O valor exato da constante é $\sqrt{2}$.
    \item \textbf{Como sabemos disso?} Como o próprio texto aponta, a análise dimensional e a consistência dimensional têm o poder de nos dizer a forma da relação entre as variáveis (neste caso, a proporcionalidade com a raiz de $gh$), mas \textbf{elas falham em determinar o valor de constantes numéricas adimensionais} [1]. Para descobrir o valor exato dessa constante, precisamos obrigatoriamente recorrer à teoria física, como as leis de Newton, ou realizar experimentos [1].
\end{itemize}

Utilizando os princípios da mecânica clássica, podemos resolver o problema apelando para a \textbf{Lei da Conservação de Energia}. Quando um corpo cai do repouso a partir de uma altura $h$, toda a sua energia potencial gravitacional inicial ($E_p$) se transforma em energia cinética ($E_k$) instantes antes de atingir a referência zero:
$$ E_p = E_k $$
$$ mgh = \frac{1}{2}mV^2 $$

Cancelando a massa ($m$) de ambos os lados e isolando a velocidade ($V$), obtemos:
$$ V^2 = 2gh \implies V = \sqrt{2} \cdot \sqrt{gh} $$

Ao compararmos este resultado analítico exato da mecânica newtoniana com a equação (2.12) obtida apenas por agrupamento dimensional, verificamos inequivocamente que a "constante" que faltava no modelo é $\sqrt{2}$.
